\section{Introduction}
\label{sec:introduction}

Industrial maintenance alarms are useful only when they map to the decision unit that operators and engineers actually face. In some settings the unit is a failure episode, in others a repair history, a laboratory component-condition cycle, or a production example with a downstream yield label. Standard pointwise anomaly detection can be accurate at the sample level while still producing a poor operational object. A burst of extreme samples near one physical failure may be counted as many alarms; a vehicle-level repair problem may be misread as timestamp-level warning; and millions of telemetry rows can obscure the small number of independent events or entities that actually support inference.

Cyclic machinery creates an additional statistical problem. Compressor, engine, and actuator signals are shaped by operating regimes, control loops, maintenance states, and sensor preprocessing. Extreme observations therefore need not be independent exceedances of a stationary distribution. Dynamical extreme-value theory connects high observations of functions evaluated along a trajectory with returns to small target sets, and it explains why periodic or recurrent structure can create clustered extremes \citep{Freitas2008Hitting,Freitas2012Periodicity}. Under finite noisy observations, the same asymptotic quantities may be masked or distorted \citep{Faranda2013Noise}. These facts make naive threshold exceedance counts an unstable basis for maintenance claims.

This paper studies a constrained contribution: a reproducible diagnostic workflow that turns machine-state trajectories and tabular condition-monitoring examples into artifact-backed evidence, with explicit warnings about what the artifacts do not prove. The workflow defines causal state vectors, operating regimes, training-defined dangerous regions, regime-conditioned thresholds, exceedance clusters, lagged multivariate patterns, alarm episodes, and dataset-specific condition diagnostics evaluated under a frozen protocol. The central question is not whether a new model is a general winner across industrial settings. The question is whether generated artifacts make dependence, clustering, alarm burden, sensitivity, calibration, estimand changes, and unsupported claims visible enough to constrain a defensible predictive-maintenance paper.

The contributions are therefore deliberately bounded. First, the implementation provides a dynamical-EVT-inspired formulation for regime-conditioned extreme diagnostics in predictive maintenance. Second, it produces finite-sample simulation and sensitivity artifacts for extremal-index recovery and threshold/run-length choices. Third, it runs a five-dataset real-data matrix covering compressor telemetry, fleet repair histories, hydraulic component-condition cycles, and semiconductor process/yield examples. Fourth, it freezes evaluation and audit protocols so recall, false alarms, duplicate alarms, lead time, calibration, verification status, and industrial diagnostic rows are generated rather than hand-curated. Finally, it generates a claim ledger that allows only artifact-backed statements in the manuscript and provides build commands that regenerate the figures, tables, provenance, manuscript, supplement, and audits from version-controlled code. These contributions correspond to claim-ledger entries CLM-001 through CLM-007; claims outside that ledger are not used as positive evidence.
